% ============================================================================
% Fichier  : partie4/P04_C12_methodes_risque.tex
% Titre    : Les Méthodes d'Appréciation des Risques
% Auteur   : MINKA MI NGUIDJOÏ Thierry Emmanuel
% ============================================================================
% Sources : EBIOS Risk Manager (ANSSI, 2018) ; OCTAVE ; MEHARI ; ISO 27005.
% Mairie d'Eséka : cas fictif.
% ============================================================================

\chapter{Les Méthodes d'Appréciation des Risques}
\label{chap:methodes}
\chaptermark{Les méthodes de risque}

\epigraph{%
    « Une méthode ne supprime pas le jugement~:\\
    elle le rend traçable et partageable. »%
}{Principe du cours}

\niveauChapitre{Intermédiaire}{Le Chapitre~\ref{chap:vocabrisque}.}

\begin{boxobjectifs}
\begin{itemize}
    \item Présenter \sigle{EBIOS} Risk Manager et ses cinq ateliers.
    \item Décrire la méthode \sigle{OCTAVE} et ses trois phases.
    \item Situer \sigle{MEHARI} et son approche par bases de connaissances.
    \item Comparer les méthodes et choisir selon le contexte.
    \item Appliquer une démarche d'appréciation à la mairie.
\end{itemize}
\end{boxobjectifs}

\bigskip

% ============================================================================
\section*{Question directrice : comment conduire une appréciation des risques sans se perdre ?}
% ============================================================================

Apprécier les risques « à l'intuition » produit des oublis et des injustices.
Les méthodes structurent la démarche. Ce chapitre présente les principales et
aide à choisir.

\medskip
\textbf{Plan du chapitre.}
\begin{enumerate}
    \item EBIOS Risk Manager (section~\ref{sec:c12-ebios}).
    \item OCTAVE (section~\ref{sec:c12-octave}).
    \item MEHARI (section~\ref{sec:c12-mehari}).
    \item Comparer et choisir (section~\ref{sec:c12-choisir}).
    \item Appliquer à la mairie (section~\ref{sec:c12-mairie}).
\end{enumerate}

\separateur

% ============================================================================
\section{EBIOS Risk Manager}
\label{sec:c12-ebios}
% ============================================================================

\begin{boxdef}
\termecle{\sigle{EBIOS} Risk Manager}~: méthode française d'appréciation et de
traitement des risques numériques, publiée par l'\sigle{ANSSI}, organisée en
cinq ateliers.
\end{boxdef}

\begin{tableaucours}{L{1.4cm}L{4.6cm}L{6.2cm}}{%
    \tableauHeader{Atelier} &
    \tableauHeader{Nom} &
    \tableauHeader{Produit}
}
1 & Cadrage et socle & Périmètre, valeurs métier, socle de sécurité \\
2 & Sources de risque & Qui menace, pour quels objectifs \\
3 & Scénarios stratégiques & Chemins d'attaque de haut niveau \\
4 & Scénarios opérationnels & Chemins techniques détaillés \\
5 & Traitement du risque & Mesures, risques résiduels \\
\end{tableaucours}
\vspace{-0.8em}
\begin{center}{\small\itshape Tableau~12.1. Les cinq ateliers d'EBIOS Risk Manager}\end{center}

\sigle{EBIOS} \sigle{RM} met l'accent sur les \emph{sources de risque} et les
\emph{scénarios}~: elle part de qui veut nuire et pourquoi, plutôt que d'une
liste de vulnérabilités.

% ============================================================================
\section{OCTAVE}
\label{sec:c12-octave}
% ============================================================================

\begin{boxdef}
\termecle{\sigle{OCTAVE}} (\emph{Operationally Critical Threat, Asset and
Vulnerability Evaluation})~: méthode centrée sur les actifs critiques,
organisée en trois phases.
\end{boxdef}

\begin{enumerate}
    \item \textbf{Profils de menaces fondés sur les actifs}~: identifier les
          actifs critiques et les menaces qui les visent.
    \item \textbf{Vulnérabilités de l'infrastructure}~: analyser les
          faiblesses techniques.
    \item \textbf{Stratégie et plan de sécurité}~: évaluer impact et
          probabilité, bâtir le plan.
\end{enumerate}

\sigle{OCTAVE} est conduite par l'organisation elle-même, ce qui favorise
l'appropriation.

% ============================================================================
\section{MEHARI}
\label{sec:c12-mehari}
% ============================================================================

\begin{boxdef}
\termecle{\sigle{MEHARI}} (Méthode Harmonisée d'Analyse de Risques)~: méthode
du \sigle{CLUSIF} s'appuyant sur des bases de connaissances pour évaluer
l'exposition et l'efficacité des mesures.
\end{boxdef}

\sigle{MEHARI} propose une analyse fine en situation normale, en cas de
dysfonctionnement majeur, et par la vulnérabilité. Son atout est la richesse de
ses bases~; son coût, leur lourdeur pour une petite structure.

% ============================================================================
\section{Comparer et Choisir}
\label{sec:c12-choisir}
% ============================================================================

\begin{tableaucours}{L{2.6cm}L{4.4cm}L{5.2cm}}{%
    \tableauHeader{Méthode} &
    \tableauHeader{Force} &
    \tableauHeader{Convient à}
}
EBIOS RM & Scénarios, sources de risque & Projets structurés, lien ISO 27005 \\
OCTAVE & Appropriation interne & Organisations qui s'auto-évaluent \\
MEHARI & Bases de connaissances & Structures matures et outillées \\
\end{tableaucours}
\vspace{-0.8em}
\begin{center}{\small\itshape Tableau~12.2. Choisir une méthode selon le contexte}\end{center}

Aucune méthode n'est « la meilleure »~: la bonne est celle que l'organisation
peut réellement conduire. Toutes sont compatibles avec \norme{ISO/IEC 27005}.

% ============================================================================
\section{Appliquer à la Mairie}
\label{sec:c12-mairie}
% ============================================================================

\begin{boxprocedure}
\textbf{Constatation.} La mairie d'Eséka, aux moyens limités, choisit
\sigle{EBIOS} \sigle{RM} pour l'opération \textsc{Kellé}, en version allégée.

\medskip
\textbf{Analyse.} Le choix est cohérent~: \sigle{EBIOS} \sigle{RM} est
publique, gratuite, outillée, et son approche par scénarios parle aux élus (qui
veut nuire à la mairie, et comment). Une version allégée évite la lourdeur pour
une petite structure.
\end{boxprocedure}

\begin{boxmetier}{RSSI}
\textbf{Adapter, pas appliquer à la lettre.} La RSSI adapte la méthode à la
taille de la mairie~: cinq ateliers menés en une journée d'atelier collectif,
plutôt qu'une étude de plusieurs mois. La méthode sert le projet, pas
l'inverse.
\end{boxmetier}

\separateur

\begin{boxresume}
\textbf{Ce qu'il faut retenir de ce chapitre~:}
\begin{itemize}
    \item \textbf{EBIOS RM} (ANSSI) structure l'appréciation en cinq ateliers,
          centrés sur les scénarios.
    \item \textbf{OCTAVE} part des actifs critiques en trois phases, par
          auto-évaluation.
    \item \textbf{MEHARI} s'appuie sur des bases de connaissances, pour
          structures matures.
    \item On \textbf{choisit} la méthode que l'organisation peut conduire~;
          toutes sont compatibles ISO 27005.
    \item La mairie adopte \textbf{EBIOS RM allégée} pour l'opération
          \textsc{Kellé}.
\end{itemize}
\end{boxresume}

\begin{boxinfo}
\textbf{Points de contrôle de l'élève.} Je sais~: citer les cinq ateliers
d'EBIOS RM~; décrire les trois phases d'OCTAVE~; situer MEHARI~; comparer les
méthodes~; justifier un choix.
\end{boxinfo}

\bigskip

\begin{boxexo}[title={\ding{45}\quad Exercice~12.1. Les ateliers \niveaubase}]
Citez les cinq ateliers d'\sigle{EBIOS} Risk Manager.
\end{boxexo}

\begin{boxexo}[title={\ding{45}\quad Exercice~12.2. Les phases \niveaubase}]
Énoncez les trois phases d'\sigle{OCTAVE}.
\end{boxexo}

\begin{boxexo}[title={\ding{45}\quad Exercice~12.3. MEHARI \niveaubase}]
Quel est le principe de \sigle{MEHARI} et son principal atout~?
\end{boxexo}

\begin{boxexo}[title={\ding{45}\quad Exercice~12.4. Source de risque \niveaumoyen}]
En quoi EBIOS RM se distingue-t-elle en partant des sources de risque~?
\end{boxexo}

\begin{boxexo}[title={\ding{45}\quad Exercice~12.5. Choisir \niveaumoyen}]
Pour une petite mairie, quelle méthode recommandez-vous~? Justifiez.
\end{boxexo}

\begin{boxexo}[title={\ding{45}\quad Exercice~12.6. Adapter \niveaumoyen}]
Pourquoi faut-il adapter la méthode à la taille de l'organisation~?
\end{boxexo}

\begin{boxexo}[title={\ding{45}\quad Exercice~12.7. Atelier 1 \niveaumoyen}]
Que produit l'atelier~1 d'EBIOS RM pour la mairie d'Eséka~?
\end{boxexo}

\begin{boxexo}[title={\ding{45}\quad Exercice~12.8. Compatibilité ISO \niveaumoyen}]
Pourquoi ces méthodes sont-elles toutes compatibles avec \norme{ISO/IEC 27005}~?
\end{boxexo}

\begin{boxexo}[title={\ding{45}\quad Exercice~12.9. Mini-EBIOS \niveauexpert}]
Conduisez un mini-atelier EBIOS RM sur l'état civil~: deux sources de risque,
deux scénarios stratégiques.
\end{boxexo}

\begin{boxexo}[title={\ding{45}\quad Exercice~12.10. Note de méthode \niveauexpert}]
Rédigez la note qui justifie, auprès du maire, le choix de méthode pour
l'opération \textsc{Kellé}.
\end{boxexo}

\bigskip

\begin{boxinfo}
\textbf{Pour aller plus loin.}
\begin{itemize}
    \item EBIOS Risk Manager~\cite{ebios2018}~; ISO/IEC 27005~\cite{iso27005}.
\end{itemize}
\end{boxinfo}

\begin{boxinfo}
\textbf{Transition.} Les risques appréciés, il faut les cartographier et
décider~: le Chapitre~\ref{chap:traitement} traite la cartographie et le
traitement.
\end{boxinfo}
